%!TEX program = xelatex
\documentclass[11pt,a4paper]{article}
\usepackage[margin=2.2cm]{geometry}
\usepackage{amsmath,amssymb,amsthm,mathtools}
\usepackage{bm}
\usepackage{enumitem}
\usepackage{booktabs}
\usepackage{xcolor}
\usepackage{hyperref}
\usepackage{fancyhdr}
\usepackage{xeCJK}

\hypersetup{
    colorlinks=true,
    linkcolor=blue!60!black,
    citecolor=blue!60!black,
    urlcolor=blue!60!black
}

\pagestyle{fancy}
\fancyhf{}
\rhead{\small Qiuzhen QE Probability \& Statistics --- Fall 2025}
\lhead{\small Official Qualifying Exam}
\rfoot{\small Page \thepage}

\DeclareMathOperator*{\argmax}{arg\,max}
\DeclareMathOperator*{\argmin}{arg\,min}
\DeclareMathOperator{\Var}{Var}
\DeclareMathOperator{\E}{\mathbb{E}}
\DeclareMathOperator{\Pbb}{\mathbb{P}}
\DeclareMathOperator{\Cov}{Cov}

\setlist[itemize]{topsep=3pt,itemsep=2pt,parsep=0pt}
\setlist[enumerate]{topsep=3pt,itemsep=2pt,parsep=0pt}

\title{\textbf{QUALIFYING EXAM: PROBABILITY \& STATISTICS}\\[0.5em]\Large FALL 2025}
\date{}
\author{}

\begin{document}

\maketitle
\vspace{-3em}

\section*{Instructions}
\begin{itemize}
    \item There are 11 problems in this exam. You need to choose 8 of them to solve. If you select more than 8, only the first 8 that you have worked on will be graded. Note that 4 of the problems are worth 15 points each and the rest 10 points each.
    \item You must follow all the rules of exam taking. Misconducts will be subject to proper disciplinary actions by the Center.
    \item You must provide all necessary details for full credits. A final answer with no or little explanation/derivation, even if correct, receives a minimal credit.
    \item $\mathbb{R}$ denotes the set of real numbers and $\mathbb{N} = \{1, 2, 3, \dots\}$ denotes the set of positive integers. $\xrightarrow{(d)}$ and $\stackrel{(d)}{=}$ mean ``converges in distribution'' and ``equal in distribution'', respectively.
    \item Beta function $\mathrm{Beta}(\alpha, \beta) = \int_0^1 x^{\alpha-1}(1-x)^{\beta-1} dx = \frac{\Gamma(\alpha)\Gamma(\beta)}{\Gamma(\alpha+\beta)}$.
    \item For $\mu \in \mathbb{R}, \sigma > 0$, $\mathcal{N}(\mu, \sigma^2)$ density is $f(x \mid \mu, \sigma^2) = \frac{1}{\sqrt{2\pi\sigma^2}} \exp\left(-\frac{(x-\mu)^2}{2\sigma^2}\right)$.
\end{itemize}

\noindent\rule{\linewidth}{0.4pt}

\begin{enumerate}[label=\textbf{\arabic*.}]
    \item \textbf{(10 points)} Let $X$ and $N$ be independent random variables. $X \sim \mathrm{Exp}(1)$ (density $e^{-x}\mathbf{1}_{\{x > 0\}}$) and $N \sim \mathrm{Poisson}(\lambda)$ for $\lambda > 0$.
    \begin{enumerate}[label=(\alph*)]
        \item Find all $\lambda > 0$ such that $\E[X^N] < \infty$.
        \item Calculate $\E[X^N]$ under the above condition on $\lambda$.
        \item Is there any $\lambda > 0$ such that $\Var[X^N] < \infty$? Justify your answer.
    \end{enumerate}

    \item \textbf{(10 points)} Let $X, Y$ be independent uniform on $[0, 1]$, and $Z \coloneqq \max\{X, Y\}$.
    \begin{enumerate}[label=(\alph*)]
        \item Determine the distribution of $Z$ conditioned on $\{X + Y \in [0, 1]\}$.
        \item Determine the distribution of $Z$ conditioned on $\{X + Y \in [1, 2]\}$.
    \end{enumerate}

    \item \textbf{(10 points)} Let $(X_n, \mathcal{F}_n)$ be a martingale.
    \begin{enumerate}[label=(\alph*)]
        \item Prove that if $(X_n^2, \mathcal{F}_n)$ is also a martingale, then $X_m = X_n$ a.s.\ for any $m, n \in \mathbb{N}$.
        \item Prove that if $(|X_n|^p, \mathcal{F}_n)$ is a martingale for some $p > 1$, then $(|X_n|^q, \mathcal{F}_n)$ is also a martingale for every $1 \le q \le p$.
    \end{enumerate}

    \item \textbf{(15 points)} Given independent random variables $(X_n)_{n \in \mathbb{N}}$ satisfying
    \[
    \forall n \in \mathbb{N}, \qquad \Pbb(X_n = -1) = \frac{1}{2}, \qquad \Pbb(X_n = 4n) = \frac{1}{8n}, \qquad \Pbb(X_n = 0) = \frac{1}{2} - \frac{1}{8n},
    \]
    and denote $S_n \coloneqq \sum_{k=1}^n X_k$. Prove that almost surely
    \[
    \liminf_{n \to \infty} \frac{S_n}{n} < 0 < \limsup_{n \to \infty} \frac{S_n}{n}.
    \]

    \item \textbf{(10 points)} Let $B_t$ be 1D standard Brownian motion with $B_0 = 0$. Define $X_t = \exp(B_t - \frac{1}{2}t)$.
    \begin{enumerate}[label=(\alph*)]
        \item Prove that $X_t$ converges almost surely as $t \to \infty$.
        \item For $t > 0$, calculate $\E \int_0^t X_s ds$ and $\Var \int_0^t X_s ds$.
        \item Show that $\int_0^\infty X_t dt < \infty$ a.s.
    \end{enumerate}

    \item \textbf{(15 points)} Let $M = \begin{pmatrix} X_1 & Y \\ Y & X_2 \end{pmatrix}$ be a symmetric random matrix, where $X_1, X_2 \stackrel{\text{i.i.d.}}{\sim} \mathcal{N}(0, 2)$ and $Y \sim \mathcal{N}(0, 1)$ are independent. Let $\lambda_1 \ge \lambda_2$ be eigenvalues of $M$, and $\mathbf{v}_1, \mathbf{v}_2$ unit eigenvectors.
    \begin{enumerate}[label=(\alph*)]
        \item Find the distribution of $\lambda_1 + \lambda_2$.
        \item Find the distribution of $\lambda_1 - \lambda_2$.
        \item Are $\lambda_1$ and $\lambda_2$ independent? Prove your claim.
        \item Prove that the joint distribution of $(\mathbf{v}_1, \mathbf{v}_2)$ is invariant under orthogonal transformations.
    \end{enumerate}

    \item \textbf{(10 points)} Suppose $\hat{\theta}_n$ is a real-valued consistent estimator of $\theta \in [-1, 1]$ based on $n$ observations. Let $T_n$ be a sufficient statistic for $\theta$. Show there exists a consistent estimator of $\theta$ based on $T_n$.

    \item \textbf{(15 points)} Let $X_1, \dots, X_n \stackrel{\text{i.i.d.}}{\sim} \mathrm{Exp}(\lambda)$ (pdf $\lambda e^{-\lambda x}$ for $x > 0$). Let $\phi = P_\lambda(X_1 > x) = e^{-\lambda x}$.
    \begin{enumerate}[label=(\alph*)]
        \item Find a complete and sufficient statistic for $\lambda$.
        \item Find the UMVUE of $\phi$.
    \end{enumerate}

    \item \textbf{(15 points)} Let $X_1, \dots, X_n \stackrel{\text{i.i.d.}}{\sim} \mathrm{Bernoulli}(p)$ with $0 < p < 1$.
    \begin{enumerate}[label=(\alph*)]
        \item Let $g(p) = p^k + (1-p)^{n-k}$ ($0 \le k \le n$). Find the UMVUE for $g(p)$.
        \item Prove that $\overline{X}$ is admissible for $p$ under loss $l(p, a) = \frac{(p-a)^2}{p(1-p)}$.
    \end{enumerate}

    \item \textbf{(10 points)} Let $X_1, \dots, X_n \stackrel{\text{i.i.d.}}{\sim} \mathrm{Uniform}(0, \theta)$ with $\theta > 0$. Prior on $\theta$ is lognormal ($\ln\theta \sim \mathcal{N}(\mu_0, \sigma_0^2)$).
    \begin{enumerate}[label=(\alph*)]
        \item Find the posterior density of $\ln\theta$.
        \item Find the MAP (mode) estimator of $\theta$. Is it consistent? Explain.
    \end{enumerate}

    \item \textbf{(10 points)} Let $X_1, \dots, X_n$ be i.i.d.\ Discrete Uniform on $\{1, 2, \dots, \theta\}$.
    \begin{enumerate}[label=(\alph*)]
        \item For testing $H_0 : \theta \le \theta_0$ vs. $H_1 : \theta > \theta_0$, show that
        \[
        \phi(\mathbf{x}) = \begin{cases} 1, & \text{if } x_{(n)} > \theta_0, \\ \alpha, & \text{if } x_{(n)} \le \theta_0 \end{cases}
        \]
        is a randomized UMP test of size $\alpha$.
        \item For testing $H_0 : \theta = \theta_0$ vs. $H_1 : \theta \neq \theta_0$, show that
        \[
        \phi(\mathbf{x}) = \begin{cases} 1, & \text{if } x_{(n)} > \theta_0 \text{ or } x_{(n)} \le \theta_0 \alpha^{1/n}, \\ 0, & \text{otherwise} \end{cases}
        \]
        is a UMP test of size $\alpha$ when $\theta_0 \alpha^{1/n}$ is an integer.
    \end{enumerate}
\end{enumerate}

\end{document}
